\begin{figure}[htbp]
\centering
\begin{tikzpicture}[>=Stealth, every node/.style={font=\small},
 space/.style={circle,draw=black!60,minimum size=9mm,inner sep=2pt},
 edge/.style={->,thick,blue!65!black}]
 \node[space] (h0) at (0,1.9) {$H_0$};
 \node[space] (h1) at (-1.8,-.8) {$H_1$};
 \node[space] (h2) at (1.8,-.8) {$H_2$};
 \draw[edge] (h1) -- node[left=5pt] {$W_0$} (h0);
 \draw[edge] (h2) -- node[below=5pt] {$W_1$} (h1);
 \draw[edge] (h0) -- node[right=5pt] {$W_2$} (h2);
 \node[below=18pt] at (0,-.8) {the cyclic product $W_0W_1W_2$};
 \node[space] (r1) at (5,1.3) {$H_1$};
 \node[space] (r2) at (5,-.8) {$H_2$};
 \node[space] (r0) at (7.5,.25) {$H_0$};
 \draw[edge] (r1) -- node[above=5pt] {$W_0$} (r0);
 \draw[edge] (r2) -- node[below=5pt] {$W_2^*$} (r0);
 \node[below=18pt] at (6.25,-.8) {$R_0=[\,W_0\quad W_2^*\,]$};
\end{tikzpicture}
\caption{The three edge maps and the row star at vertex $0$.
Taking the adjoint of $W_2$ makes both blocks of $R_0$
have target $H_0$. The two Grams $R_0^*R_0$ and $R_0R_0^*$
have the same positive eigenvalues on different spaces.}
\label{fig:triangle}
\end{figure}
